\section{Related Work}\label{sec:related}

Both byte-pair encoding~\cite{sennrich2016bpe} and the unigram language model~\cite{kudo2018subword}, popularized in the SentencePiece implementation~\cite{kudo2018sentencepiece}, are subword tokenizers that learn token boundaries only from the statistics of characters.
\citet{galle2019bpe} attributes the success of BPE to the shorter sequences it produces.
\citet{mielke2021between} discuss the space of choices between characters and words, and \citet{bostrom2020bpe} observed that, in English, unigram segmentations are closer to morphemes and that language models trained on them were at least as good.
We reproduce that finding here for five more languages.

Unsupervised morphology induction dates back decades before the advent of neural language models (Morfessor~\cite{creutz2007morfessor} being a prime example) and more recently there has been effort to leverage morphological information in BPE or WordPiece vocabularies.
MorphPiece~\cite{jabbar2023morphpiece} injects morphological information by using morphological dictionaries with BPE, MorphBPE~\cite{asgari2025morphbpe} constrains BPE merges based on morpheme boundaries produced by a morphological analyzer, MorphTok~\cite{brahma2025morphtok} applies morphological segmentation to Indic text before subword segmentation, and KinyaBERT~\cite{nzeyimana2022kinyabert} constructs a morphology-aware encoder for Kinyarwanda.
\citet{hofmann2021superbizarre} show that segmenting words at derivational boundaries improves BERT's comprehension of complex words, and \citet{hofmann2022flota} obtain a similar effect by choosing the segmentation with the lowest number of tokens from a fixed vocabulary.
Lemma2Word falls into this category, though instead of choosing a morphological segmentation from a dictionary, we derive boundaries from a lemmatizer.

In the case of unspaced languages it has been common practice to perform segmentation and POS tagging jointly rather than sequentially, with improvements in performance~\cite{zhang2008joint}.
We are not aware of any prior work on subword tokenization that makes use of POS information to weight merge operations, and the POS-weighted BPE method evaluated herein is new in that sense, which we use to assess what this methodology is capable of, as well as what it is not.

Evaluation of morphological alignment with word segmentation lexica uses precision and recall metrics~\cite{arnett2025morphscore,batsuren2024evaluating}, and weak alignment has been linked to poorer language modeling performance for morphologically complex languages~\cite{arnett2025why}.
The gold standard referenced here is MorphScore~\cite{arnett2025morphscore}, which is computed from word boundaries derived using Universal Dependencies lemmas for 70 languages.
Intrinsic metrics like compression and R\'enyi efficiency~\cite{zouhar2023noiseless} are imperfect proxies for performance in downstream tasks.
As observed by \citet{schmidt2024tokenization}, the number of tokens is not a very good metric of progress.
\citet{ali2024tokenizer} report that tokenizer choice is measurable at the 2.6B model scale.
\citet{rust2021tokenizer} found that monolingual tokenizers beat a single shared multilingual tokenizer.
\citet{park2021morphology} used morphology to predict language-modeling difficulty on 92 languages.

\citet{toraman2023tokenization} evaluated tokenizers specifically for Turkish, one of the languages considered here; they find that morphology-level tokenization is useful for small vocabularies.
Our language-model probe follows the practice of these studies, but on a much more modest scale.
